\renewcommand{\stagefigureprefix}{3}\setcounter{figure}{0}

\reporttitle{第三阶段：协同仿真与性能分析}

本阶段完成NEMU指令扩展与RTL协同仿真，使用现有RVV指令构造等价对照，并评估点积及矩阵向量计算的周期与实现成本。最终硬件版本记为0f，包含两槽弹性执行结构（O2）、低32位结果保存（H1）和DIV取消关钟修复。软件中的firstdot表示首段点积初始化方法，C8表示8×32调用中的向量配置复用入口。\par

\section{NEMU扩展与协同仿真}

\subsection{独立参考模型与架构行为}

VDot的硬件实现和NEMU参考模型共享同一份指令语义。输入是两个向量寄存器低64位中的八对有符号字节，输出是八项乘积之和：低32位保存二补码结果，高96位清零。硬件通过并行乘法、归约和两槽控制产生响应，NEMU则按软件顺序提取元素并求和。两种实现采用不同的计算组织，使比较同时覆盖算术结果和结果布局。\par


参考模型的取数过程从架构寄存器开始。每个字节先解释为signed int8，再扩展到能够容纳乘积与总和的有符号类型；求和初值为零，不读取目的寄存器的旧值。读取两个源的快照后再写回目的寄存器，因此vd与vs1、vs2重合时仍使用同一条指令执行前的源数据。这样的步骤与软件可见语义直接对应，也便于检查C语言整数提升、符号扩展和移位处理是否符合预期。\par


参考侧使用完整128位向量值承载结果。对负数而言，符号只保存在低32位二补码中，高96位仍为零；这是本指令约定的结果格式。若把32位结果符号扩展到整个向量，低位算术可能正确，高位比较仍会立即暴露布局错误。硬件内部采用32位结果寄存器后，在公开输出处补零，正是利用了这一固定的架构属性。\par

\reportfigure{S3-01}{RTL与NEMU独立模型的DiffTest框架}{12.3cm}{相同程序与初始状态分别进入RTL和NEMU。左侧八路乘法和组合归约、右侧C循环仅表示独立求值结构，树层数不是流水级数，循环次数不是DUT周期。算术结果写入各自向量状态，CSR与异常来自模型整体；DiffTest按对应架构提交比较完整结果及状态。D、R仅为两端结果标签，FU响应不等于ROB提交。}{精确SVG工程图；工程依据：冻结EMU/NEMU身份、差分协议和R5结果；原图存档与所选资产身份见\seqsplit{figures/redesigned-20261105/selected-diagrams.json}。稳定0f＋固定reference；同一镜像}

\subsection{指令译码与异常处理}

NEMU先按custom-0指令字匹配项目编码，再检查向量配置。正常路径使用SEW8、LMUL1、VL8、vstart0和VLEN128；非法配置进入异常处理，目的向量保持原值。译码与执行检查共同决定指令能否生效，因此测试既覆盖操作码和功能字段，也覆盖同一合法机器字在不同向量状态下的行为。\par


香山固定版本中，译码阶段识别的非法指令与功能单元发现的VL非法具有不同的xTVAL表现。平台适配按对应的指令和架构状态描述这些条件，并比较异常码、陷阱PC、xTVAL及目的状态。报告中的特权执行记录来自M模式。该平台规则集中保存在reference配置中，正常算术路径仍由独立求和得到结果。\par


异常测试中的初始状态与正常测试同样重要。将目的向量设置为可辨认的非零值，能观察非法操作是否错误清零或部分更新目的；源寄存器和CSR快照则用于核对异常前后的状态变化。把这些观察与异常PC对应起来，可以区分译码错误、执行阶段非法配置和异常后的状态恢复问题。\par

\reportfigure{S3-02}{NEMU中vdot的识别、状态检查与独立执行}{16.9cm}{精确匹配vdot编码后，所有向量状态条件均须满足；vtype合法包括vill与保留位检查。合法路径先快照两源，再求八项有符号INT8乘积之和，完整写回128位（高96位为零、低32位为和的补码）并将VS置脏。其他译码与状态非法分支分别处理。NEMU a48a11f、补丁1.1.0；xTVAL沿用固定平台策略，已验证范围为M模式。}{精确SVG工程图；工程依据：\seqsplit{team/patches/nemu}及冻结reference实现；原图存档与所选资产身份见\seqsplit{figures/redesigned-20261105/selected-diagrams.json}。按已实现版本绘制；平台xTVAL范围按正文保留}

\subsection{提交边界上的差分比较}

DiffTest在架构提交边界比较DUT与参考模型。前端取指、功能单元响应和ROB提交分别代表不同的进度，比较器以最终提交的指令与架构状态建立对应关系。对VDot而言，关键内容包括机器指令、目的寄存器、完整向量结果及向量CSR；对异常指令，还包括陷阱入口和相关异常字段。\par


运行记录按启动、指令执行、差分检查、程序结束四个阶段组织。正常用例同时检查计划中的测试数量、完整结果、GOOD TRAP和进程退出状态。这样可以把程序自行报告的功能结果与外部差分检查联系起来：前者核对测试期望，后者核对DUT提交状态与独立模型的一致性。保留stdout、stderr和退出码，也便于定位构建、装载和目标程序执行之间的故障位置。\par


不同层次的测试采用各自的计数单位。主机模型以输入向量和测试方法计数，局部RTL以事务与控制场景计数，全核以程序、用例和窗口计数。R5全核功能矩阵的117个不同用例与124次分片执行分别描述用例覆盖和实际执行安排；新控制诊断中的131072字节则描述被核对的寄存器结果量。把这些单位列清楚，能够准确反映每一层验证所完成的工作。\par

点积测试检查整个目的向量，同时记录其他向量寄存器与CSR。VDot路径仅更新约定的目的寄存器；RVV等价路径需要中间向量寄存器，检查器据预先声明的scratch集合判断这些变化。由此既能检查计算结果，又能区分允许的临时状态与意外破坏。\par


源与目的别名是测试的重要部分。先读取两个源、再清理或覆盖目的，能够保证vd与任一源重合时仍计算原来的八对输入。RVV序列也遵循相同的数据顺序：先生成宽化乘积，再清零目的并归约。若先清目的寄存器，别名场景中的源数据会被破坏，即使无别名样例全部通过也会留下语义缺口。\par


访存测试还检查输出槽和前后护栏。每个结果使用16字节槽保存完整向量，连续结果之间的步长与实际存储宽度一致。逐槽比较能够发现结果正确但写入地址错误、前一槽覆盖后一槽或高位残留等问题。输入高半区的非零哨兵则检查实现是否意外使用了规定输入范围之外的数据。\par

\subsection{高位注错与正常参考对照}

bit100负控针对向量结果的高位比较。合法VDot结果的高96位均为零，选择完整向量的bit100注错，相当于改变高64位中的bit36。低32位算术结果可以保持一致，而高位差异应被DiffTest检出，因此这项测试直接检查完整向量比较是否生效。\par


图中的两组记录来自2026年10月24日core-01实验。负控运行加载故障注入参考库，在PC 0x80001bdc、指令字0x0294050b、目的寄存器v10处出现首次失配，以exit 1结束；正常对照使用同一测试镜像和正常reference，完成15个用例并以GOOD TRAP、exit 0结束。两次独立运行的镜像一致，使参考库成为对照中的关键变化项。\par


首次架构失配PC定位状态差异的起点，终止事件记录模拟器的退出位置。机器字、寄存器和高位差异共同对应注错目标；正常对照确认同一镜像、启动流程和数据能够完成整组测试。\par

\reportfigure{S3-03}{bit100负控与正常参考对照}{13.2cm}{负控在v10的bit100检出差异，首次失配PC为0x80001bdc；正常参考对照完成15个用例并以GOOD TRAP、exit 0结束。}{e2c2e8d固定归档：\seqsplit{team/optimizations/o2/r3-div-gate/fullcore/evidence/20261024-fullcore/files/runs/core-01/}\{negative-vdot,functional-0\}/\{stdout.log,result.json\}；图数据附件保存原文、Git blob和SHA256。历史EMU 2d1db77c…与BIN b8b7e1a7…；不是R5 clean02/acf672…日志，不添加本组独立replay，不把ABORT事件PC等同首次失配PC；仅图示原记录，无新增仿真}

\section{RVV对照与周期测量方法}

\subsection{等价RVV指令序列}

两条路径均将低8个signed int8元素分别相乘并求和，目的向量低32位保存二补码结果，高96位清零。检查覆盖完整128位及寄存器副作用：vdot保持其他向量寄存器，RVV的中间操作使用预先声明的scratch集合。\par


{\small\setlength{\tabcolsep}{3pt}\begin{longtable}{>{\raggedright\arraybackslash}p{0.2\linewidth}>{\raggedright\arraybackslash}p{0.24\linewidth}>{\raggedright\arraybackslash}p{0.24\linewidth}>{\raggedright\arraybackslash}p{0.22\linewidth}}
\caption{等价RVV指令序列}\label{tab:3-1}\\
\toprule
\textbf{RVV 路径} & \textbf{点积主体指令数} & \textbf{向量 scratch} & \textbf{入口与出口}\\\midrule\endfirsthead
\caption[]{等价RVV指令序列（续）}\\
\toprule
\textbf{RVV 路径} & \textbf{点积主体指令数} & \textbf{向量 scratch} & \textbf{入口与出口}\\\midrule\endhead
{已知配置单次} & {5} & {v2、v3} & {e8/m1/VL8、tu/mu；退出恢复相同配置}\\
{通用配置单次} & {7} & {v2、v3，另用 t3/t4 保存 CSR} & {保留原 VL/vtype，测试四种 tail/mask policy}\\
{独立四路批处理} & {14 / 4 个点积} & {v16–v23} & {e8/m1/VL8、tu/mu；退出恢复相同配置}\\
\bottomrule\end{longtable}}


已知配置序列：vwmul.vv \allowbreak{}→\allowbreak{} 设置 e16/m1/VL8 \allowbreak{}→\allowbreak{} vmv.v.i vd,0 \allowbreak{}→\allowbreak{} vwredsum.vs \allowbreak{}→\allowbreak{} 恢复 e8/m1/VL8。e16 下 8 个元素正好清零整个目的寄存器；宽化归约使用 tu，保留已经清零的高位。序列先乘后清目的寄存器，从而支持源与目的别名。通用序列增加 csrr vl/vtype，末尾用 vsetvl 恢复。四路批处理共享两次vtype设置，充分利用批量调用的配置复用。\par


每次调用保存并比较全部 32 个向量寄存器，按声明排除 RVV scratch；检查 VL、vtype、vstart、vxrm、vxsat、vcsr、实际循环次数以及缓冲区前后 canary。访存组逐个比较全部 128 位输出，并检查未使用输出空间。C 判定器使用真实源文件进行原生及 UBSan 位翻转测试。\par

\subsection{计数窗口与实际执行边界}

固定版本的\seqsplit{DecodeUnit.scala}为普通CSR设置\seqsplit{waitForward/blockBackward}，而Rename.scala会解除大部分只读CSR的这些限制。初始观察中，cycle与instret指令在ROB仍有旧指令时进入，相关标志为零。计时实验因此通过实际分派、采样和退休记录确定窗口顺序。\par


软件测量边界加入fence rw,rw。该香山版本的fence通过等待旧指令、约束后续分派和flush形成边界，日志记录其与CSR及主体指令的执行关系。窗口定义依赖这套固定微架构实现，生产核的执行规则保持原样。\par


窗口结构如下（所有实验分支相同）：\par


\begin{Verbatim}[breaklines,breakanywhere,breaksymbolleft={},breaksymbolright={},fontsize=\footnotesize,baselinestretch=1]
fence; START; fence
read instret_begin; fence
read cycle_begin; fence
执行显式汇编主体（含循环控制）
fence; read cycle_end; fence
read instret_end; fence
STOP; fence
保存结果；完整状态检查；打印
\end{Verbatim}


观察器由\seqsplit{VDOT\_T13\_TRACE=1}启用，记录CSR实际采样、计数器入ROB及窗口内退休信息，并与T12观察器一起构建。它提供只读事件观察，数据通路、复位、忙闲接口和NEMU计时维持原实现。\par


窗口检查包含以下六项：\par


1. 精确匹配四条计数器指令的 PC、机器码、ROB 身份与顺序。\par
2. CSR 读值与对应退休写回数据相等，采样位于入队和退休之间。\par
3. 每条计数器的前一 fence 已退休，后一 fence 在计数器退休后才退休。\par
4. cycle 差值等于硬件全核采样 tick 差值。\par
5. instret 差值等于真实退休计数，也等于从 ELF 机器码展开有限循环得到的指令数。\par
6. 完整窗口及 UART 记录数量、顺序与 manifest 一致；真实 exit0、DiffTest 启用、完整成功 summary 后 GOOD TRAP。\par


绝对计数器可能大于2\textasciicircum{}53，解析和差值计算全程使用Python整数。FU子时钟受门控控制，局部事件先通过T12记录的父级VFEX4 accept逐事务关联，再与ROB事件比较；全核窗口周期使用统一的全局计数。\par


校准场景包括空窗口、8条整数指令和32条依赖除法。正式结果保留完整的屏障、取指和循环开销；空窗口用于观察固定成本，小工作量场景则结合这些成本解释周期。\par

\reportfigure{S3-04}{计时窗口CSR与退休边界}{15.3cm}{四条CSR分别记录cycle与instret起止值；图中标出采样和退休位置。固定0f的\seqsplit{after/dataset0/replay0/seq4}正式样本repeat=2，cycle差为71、instret差为12。周期含隔离与循环开销，不扣空窗口；START/STOP标记跨度为206个全核tick。}{固定e2c2e8d5归档的\seqsplit{20261025-performance-cost} \seqsplit{after/performance-d0-s00-r0}日志、window-result及runtime；完整路径、原文片段、Git blob与SHA256见\seqsplit{figures/evidence-20261105/csr-window-source.json}。按实际审核过的窗口绘制；不以局部FU周期替代全核计时}

\subsection{工作负载、输入与重复测量}

{\small\setlength{\tabcolsep}{3pt}\begin{longtable}{>{\raggedright\arraybackslash}p{0.2\linewidth}>{\raggedright\arraybackslash}p{0.24\linewidth}>{\raggedright\arraybackslash}p{0.24\linewidth}>{\raggedright\arraybackslash}p{0.22\linewidth}}
\caption{工作负载、输入与重复测量}\label{tab:3-2}\\
\toprule
\textbf{场景} & \textbf{n} & \textbf{点积数} & \textbf{含义}\\\midrule\endfirsthead
\caption[]{工作负载、输入与重复测量（续）}\\
\toprule
\textbf{场景} & \textbf{n} & \textbf{点积数} & \textbf{含义}\\\midrule\endhead
{single} & {1} & {1} & {已知配置完整调用窗口}\\
{chain} & {16、64、256} & {n} & {目的向量作为下一次源，观察依赖限制}\\
{batch} & {1、8、32} & {4n} & {每轮四个独立目的，比较共享配置的 RVV 批处理}\\
{stream} & {1、16、64} & {n} & {两个 128 位加载、点积、128 位存储以及指针更新}\\
{generic} & {1、16、64} & {n} & {原 vtype=192 的 CSR 保留对照}\\
\bottomrule\end{longtable}}


链式实验将上一结果的完整向量作为下一次源，数值可能逐步趋零，寄存器依赖关系则持续存在。批处理使用四个独立目的并共享输入和配置，用于观察独立请求吞吐。两者分别突出反馈等待与连续接收，对应不同的软件组织方式。\par


三组输入包括正极值、负极值和混合符号；高半区设置非零哨兵，访存组加入固定种子的随机数据。功能矩阵使用原始 22 组定向和前 32 组随机数据，轮换七种别名与四种 policy；54对数据覆盖选定的别名和策略组合。另用固定 NEMU 对全部 10,022 组数据重新检查精简基线，明确与全核执行层分开计数。\par


每个配置在同一镜像内运行6轮配对，按执行轮次保留首次样本、预热和4轮正式样本；顺序按AB/BA交替，缓存状态随实际运行演进。每种配置单独分片，两个实现的入口均 64 字节对齐；代码长度、地址和寄存器压力的自然差异保留并记录。\par


每个性能镜像以seed=1在新的EMU进程中运行三次，按相同seq比较周期和退休指令数，完成确定性重放核查。重放记录与正式统计样本分别保存。原测量批次按最多8个宿主进程安排，各DUT均采用相同单核和EMU\_THREADS=1配置；加速比使用DUT周期，宿主秒数用于记录执行成本。\par

\subsection{加速比与统计方法}

首次、预热和正式窗口完整保留。每个配置与数据集分别汇总原始周期、周期/点积、退休指令数及正式样本min/median/max；配对加速比为RVV cycles / vdot cycles，周期减少比例为1 - vdot / RVV。配对保持场景、n和数据集一致，同时列出加速和退化结果。\par


结果分为局部FU响应、完整序列窗口、独立批量吞吐和包含访存的循环四个层次。历史单在途FU的E3/II4描述局部接口行为；完整序列窗口则包含相同的隔离与循环开销。将各层指标对应到实际工作量，可以分别解释硬件执行和软件调用成本。\par


本节使用周期指标比较固定工作负载的执行效率。面积与静态时序由后续成本实验独立测量，C8软件收益则来自固定GEMV调用矩阵。三类结果分别呈现硬件周期、物理实现成本和应用调用效果。\par

正式加速比使用两种方法的总周期之比。设第 $i$ 个配对窗口的RVV周期为 $A_i$、VDot周期为 $B_i$，则汇总结果为 $\frac{\sum_i A_i}{\sum_i B_i}$。它等于各窗口加速比 $\frac{A_i}{B_i}$ 按 $B_i$ 加权后的平均值。这样的定义对应完成整组固定工作量时两种实现各自付出的总周期，较长的窗口按其实际执行成本获得相应权重。\par


逐窗加速比的简单平均回答的是另一种问题：每个窗口无论长短均占相同权重。两者在样本周期不同的情况下可能给出不同数值，因此本报告沿用原测量合同的总周期比，并在样本文件中保留每一对原值。三组数据、每组四次正式窗口的数量一致，配对身份则保证统计中的工作量一一对应。\par


对于0f修复前后比较，96对正式窗口的每对差值均为零。这比两组均值在显示精度内相同提供了更具体的信息：相同镜像和序号的窗口逐项保持周期，同时退休数也一致。面积与时序发生变化时，仍可据此把所观察到的代价定位到综合映射与路径，而不是该组软件窗口的执行周期。\par

\section{点积与矩阵向量计算性能}

\subsection{四类点积微基准}

下表来自0f同镜像对照合同，均为正式窗口均值。\par


{\small\setlength{\tabcolsep}{3pt}\begin{longtable}{>{\raggedright\arraybackslash}p{0.25\linewidth}>{\raggedright\arraybackslash}p{0.21\linewidth}>{\raggedright\arraybackslash}p{0.21\linewidth}>{\raggedright\arraybackslash}p{0.23\linewidth}}
\caption{四类点积微基准}\label{tab:3-3}\\
\toprule
\textbf{场景} & \textbf{vdot cycles} & \textbf{RVV cycles} & \textbf{加速比}\\ \midrule\endfirsthead
\caption[]{四类点积微基准（续）}\\
\toprule
\textbf{场景} & \textbf{vdot cycles} & \textbf{RVV cycles} & \textbf{加速比}\\ \midrule\endhead
单次点积 & 71.000 & 85.250 & 1.200704×\\
依赖链256 & 1857.000 & 5440.500 & 2.929725×\\
四路独立batch128 & 221.000 & 678.000 & 3.067873×\\
加载/计算/存储stream64 & 1127.500 & 1133.000 & 1.004878×\\
\bottomrule\end{longtable}}


独立批量可以利用O2的连续接收能力，依赖链受到前后数据依赖约束，stream场景还承担访存及循环开销。结果显示收益随工作负载改变。\par


O1→O2的独立同镜像比较中，batch128由362.5降至221周期，减少39.0345\%。H1及0f修复分别保持各自96对正式窗口周期相同。阶段间比较采用各自冻结合同。\par

\reportfigure{S3-05}{四类微基准的相对周期对照}{13.3cm}{取固定after记录，各场景、各方法均为12个正式样本。各场景按自身RVV平均窗口周期归一化为100\%，条内为实际平均窗口cycles，条外为VDot周期占比，越短越好。single、chain、batch、stream每窗口分别包含1、256、128、64次点积；batch为四路独立任务。}{冻结after summary.csv与\seqsplit{figure-data.json}交叉核对；精确Matplotlib/SVG图及固定轮廓PDF见\seqsplit{redesigned-20261105/formal/S3-05/selection.json}。局部微基准；周期占比不等于加速比，不能作为完整应用收益。加速比按总周期之比计算，见正文。}

\subsection{单次调用与独立批量}

单次点积中，VDot为71周期，RVV为85.25周期，加速比约1.201。该窗口的点积工作量较小，配置、循环和测量边界占有明显比例。VDot把多条向量算术操作收敛为一条自定义操作，节省的指令工作与共同固定开销一起形成了这一结果。\par


四路独立batch128的VDot均值为221周期，RVV为678周期，加速比约3.068。四个轮换目的为连续请求提供独立性，P与R两槽可以在同一边沿完成旧R响应、旧P迁移和新请求接收。窗口中的128次点积与循环控制共同构成总工作量，持续接收能力在这一场景中得到更充分的利用。\par


O1到O2的同镜像比较进一步隔离了流水组织的影响：batch128由362.5降到221周期。O2保留两槽状态并实现弹性推进，使前后事务能够重叠；当下游ready持续有效时，独立请求可逐拍接收。这个机制与实测周期变化相互对应，说明了优化为何对独立批量有效。\par

\subsection{依赖链与流式循环}

chain256的VDot为1857周期，RVV为5440.5周期，加速比约2.930。上一条结果写回后，下一条才能获得所需源数据，因而该场景受到完整反馈路径的约束。发射等待、寄存器读取、功能单元执行和写回共同决定下一次请求何时可发出，局部两拍响应只是其中一段。\par


stream64的VDot为1127.5周期，RVV为1133周期，加速比约1.005。两者差异较小，说明在这个加载、计算、存储和循环控制共同组成的窗口里，点积主体的优化只转化为有限的整窗收益。现有记录给出总体周期；进一步归因时，可以沿访存等待、数据依赖和单元占用逐项增加观察。\par


这四类结果共同描述VDot的使用特点：单次调用体现固定开销，独立批量体现连续接收，依赖链体现结果反馈，stream体现数据搬运与计算的组合。应用选择可以据此关注自身工作组成，例如能否把多个点积组织为独立批量、配置是否能够跨调用共享，以及结果是否立即成为下一次输入。\par

\subsection{每次点积的平均成本}

在工作量明确时，总周期还可以换算为每次点积的窗口平均成本。single、chain256、batch128和stream64分别完成1、256、128和64次点积，因而可以直接用对应均值除以点积数。这个指标把不同长度的循环放到相同工作单位下，便于观察批量组织和依赖关系的影响。\par


{\small\setlength{\tabcolsep}{3pt}\begin{longtable}{>{\raggedright\arraybackslash}p{0.2\linewidth}>{\raggedright\arraybackslash}p{0.24\linewidth}>{\raggedright\arraybackslash}p{0.24\linewidth}>{\raggedright\arraybackslash}p{0.22\linewidth}}
\caption{每次点积的平均成本}\label{tab:3-4}\\
\toprule
\textbf{场景} & \textbf{点积次数} & \textbf{VDot窗口周期/点积} & \textbf{RVV窗口周期/点积}\\\midrule\endfirsthead
\caption[]{每次点积的平均成本（续）}\\
\toprule
\textbf{场景} & \textbf{点积次数} & \textbf{VDot窗口周期/点积} & \textbf{RVV窗口周期/点积}\\\midrule\endhead
{single} & {1} & {71.0000} & {85.2500}\\
{chain256} & {256} & {7.2539} & {21.2520}\\
{batch128} & {128} & {1.7266} & {5.2969}\\
{stream64} & {64} & {17.6172} & {17.7031}\\
\bottomrule\end{longtable}}


表中数值由前述窗口均值计算，保留四位小数。batch128的1.7266周期/点积包含循环和计时边界，与功能单元逐拍接收的II=1共同描述系统的两个层次：接口可以每拍接受独立请求，软件窗口仍需要完成发射、循环及前后准备。两者之间的差异为进一步分析提供入口，也说明提高功能单元吞吐后还需考虑整条执行路径。\par


chain与batch使用不同的依赖组织。chain每次等待上一结果，batch让四个目的独立，因此归一化后的差距体现完整反馈与连续工作方式的差异。stream额外搬运输入和输出，其每次点积成本也包含访存与寻址。对实际软件进行选择时，先判断工作属于哪一种组织方式，再使用相应指标作为参考。\par

\subsection{GEMV调用与配置复用收益}

firstdot以首段点积初始化累加状态，后续段按约定求和。C8在已验证8×32路径中复用配置，提供普通C显式入口和RVV回退。输入为对齐的288字节A||x，输出为8个16字节完整槽。\par


{\small\setlength{\tabcolsep}{3pt}\begin{longtable}{>{\raggedright\arraybackslash}p{0.25\linewidth}>{\raggedright\arraybackslash}p{0.69\linewidth}}
\caption{GEMV调用与配置复用收益}\label{tab:3-5}\\
\toprule
\textbf{同合同8×32实现} & \textbf{完整调用cycles}\\ \midrule\endfirsthead
\caption[]{GEMV调用与配置复用收益（续）}\\
\toprule
\textbf{同合同8×32实现} & \textbf{完整调用cycles}\\ \midrule\endhead
firstdot & 282.25\\
原C8配置复用vdot & 273.00\\
本批最快合格RVV & 299.00\\
\bottomrule\end{longtable}}


C8相对firstdot减少3.2772\%，相对同批最快RVV加速1.095238倍。三数据集完整矩阵给出周期比较，普通C调用、软件包与最终镜像检查验证了ABI、回退和接入方式。\par


原合格镜像、复制对照、SDK接入版BIN均为31440字节，SHA相同。PT\_LOAD、48条调用表项与12组输入数据一致，ELF差异位于调试节。627项输入与封存材料的身份核对，将SDK接入结果与原功能记录对应起来。\par


4×8和16×64采用原firstdot入口，4×13选择更快的RVV。各尺寸沿用对应的已验证调用入口。\par

\subsection{硬件与软件收益的对应关系}

O2、H1和0f分别解决不同问题。O2调整两槽弹性控制，提高独立请求的重叠程度；H1利用高96位恒零的语义，将结果状态压缩为低32位保存；0f处理DIV取消后的关钟清理。把这三类改动分开，有助于分别观察吞吐、状态成本和父级控制正确性。\par


0f修复前后的192个正式窗口组成96对，逐对周期差均为零，退休指令与工作量一致。四次额外batch重放也与对应记录相同。因此该修复在这些微基准中保持原有性能，同时解决取消后持续产生无效子时钟脉冲的问题。其面积和时序影响另由完整VFEX4父级综合结果评估。\par


firstdot与C8属于软件调用层的优化。它们使用既定VDot语义，通过累加初始化、配置复用和调用组织减少软件开销。软件收益来自完整调用窗口的比较，硬件收益来自相应硬件版本的同镜像比较；这两条分析线最后汇合到选定的工程组合0f＋firstdot V2＋原C8。\par

\section{普通C接口与软件包接入}

\subsection{firstdot与C8的调用组织}

firstdot用首段点积初始化累加值，后续段按约定归约，避免在首段之前额外构造等价的初始化工作。C8针对8×32路径复用向量配置，把已验证的汇编函数接到普通C入口。两者均保留完整输出格式和相应的RVV对照。\par


在同批8×32完整调用窗口中，firstdot为282.25周期，C8为273周期，减少3.2772\%；最快合格RVV为299周期，C8相对其为1.095238倍。这个比较同时包含函数执行和该合同规定的调用成本，反映了配置复用在固定GEMV尺寸上的效果。\par


不同尺寸采用各自的已验证入口。4×8和16×64保持原firstdot路径，4×13选用较快的RVV，8×32使用原C8。这样的选择利用了实测差异：自定义点积与配置复用对某些形状有效，而较小或不规则的工作量可能更适合现有RVV序列。最终软件入口据此保留明确的选择与回退方式。\par

\reportfigure{S3-06}{C8公共调用、编译选择与输入输出布局}{14.6cm}{xs\_c8\_gemv\_8x32在编译时绑定VDot或RVV实现。packed\_ax包含行优先A[8][32]的256字节及其后的x[32]，共288字节；output\_slots包含8个16字节槽，每槽高96位清零、低32位为点积补码。两缓冲区各自16字节对齐且互不重叠。横向结果位域为高位到低位示意，不表示字节端序；VLEN=128及向量状态要求见正文。}{精确SVG工程图；工程依据：C8 CONTRACT、relink RESULT和对应冻结manifest；原图存档与所选资产身份见\seqsplit{figures/redesigned-20261105/selected-diagrams.json}。原C8稳定入口；不引入split-input或分段归约候选}

\subsection{输入输出布局与调用约定}

函数原型为 void entry(const int8\_t *packed\_ax, void *output\_slots)。已验证运行环境为当前香山昆明湖 V2/Nexus-AM 裸机 RV64、lp64d，\textbf{VLEN 恰好 128（vlenb=16）}，VS 已启用。默认路径要求该硬件实现项目的自定义 vdot 编码。调用环境须启用VS，并确认目标硬件支持项目自定义指令；本次验证采用Nexus-AM裸机环境。\par


\begin{itemize}
\item 输入：有符号 INT8 行优先 A[8][32]，紧接 x[32]，共 288 字节。
\item 输出：8 个 16 字节槽，共 128 字节。每槽低 32 位为 INT32 点积结果，高 96 位清零。槽步长为16字节，与完整向量存储一致。
\item 两块缓冲区各自 16 字节对齐、有效且不重叠；输入不会被修改。调用者负责生命周期、数据打包及平台初始化。
\item 沿用原 RV64 C ABI：栈及被调用者保存寄存器约束不变；向量寄存器和允许的调用者保存状态会改变。
\item 返回状态为 vl=8、vtype=0（e8/m1/tu/mu）、vstart=0、vcsr=7，不保留原向量配置。
\end{itemize}


底层vs1r.v存储整个向量寄存器，输出行步长固定为16字节，因此接口要求VLEN恰好128。VLEN256的整寄存器存储会超出当前槽宽。强制RVV改变调用目标，完整静态库仍保留各函数，平台与缓冲区约定保持相同。\par

\subsection{编译期选择与RVV回退}

默认的xs\_c8\_gemv\_8x32入口通过编译期别名连接配置复用VDot函数。设置\seqsplit{XS\_C8\_GEMV\_FORCE\_RVV=1}后连接相应RVV函数，显式\seqsplit{xs\_c8\_gemv\_8x32\_rvv}入口则始终选择RVV。开关只接受字面量0或1，错误配置在编译期被拒绝。\par


这一选择由头文件和匹配静态库共同形成。宏或API版本变化后，相关翻译单元一起重编，使调用点与库的选择一致。撤销C8时恢复原firstdot头文件和对应库即可回到原入口。编译期选择适用于当前明确的平台配置，调用方根据目标平台确认指令支持情况。\par

\subsection{对象、装载镜像与调用表核对}

比较对象依次为原合格C8镜像、原源码复制对照和SDK接入版。\par


{\small\setlength{\tabcolsep}{3pt}\begin{longtable}{>{\raggedright\arraybackslash}p{0.29\linewidth}>{\raggedright\arraybackslash}p{0.65\linewidth}}
\caption{对象、装载镜像与调用表核对}\label{tab:3-6}\\
\toprule
\textbf{项目} & \textbf{结果}\\\midrule\endfirsthead
\caption[]{对象、装载镜像与调用表核对（续）}\\
\toprule
\textbf{项目} & \textbf{结果}\\\midrule\endhead
{完整BIN} & {三者全字节一致，均31,440字节}\\
{BIN SHA256} & {\seqsplit{56805506eb98b68269c73a3645d8f75dff73b697e30dec5f93ccc4716089ad34}}\\
{入口} & {均0x80000000；ABI flags及ISA属性一致}\\
{PT\_LOAD} & {同地址、权限、对齐及完整内容；file30,024B，memory31,440B，零填充1,416B}\\
{装载节} & {所有SHF\_ALLOC节的位置、大小、权限、对齐及内容一致，包括BSS}\\
{函数} & {16个公共函数及启动/主程序等5个函数的地址、大小和机器码一致}\\
{测试表与数据} & {48条64字节调用记录与原计划逐项匹配；12组输入数组地址及内容一致}\\
{GNU链接} & {无未解析符号；map与实际子命令证明使用SDK库成员}\\
{完整ELF原字节} & {不同；差异仅在.debug\_info、.debug\_line、.debug\_line\_str}\\
{独立GNU复核} & {分别objcopy -{}-strip-debug后，三份完整ELF全字节一致}\\
{去调试ELF SHA256} & {\seqsplit{7e26e295fb020a63334026f768ed4825372266465b4e6008973de966fefc04ab}}\\
\bottomrule\end{longtable}}


调试节保存源文件/头文件路径及行号等信息，构建目录与头文件发生变化会改变这些记录。主检查器首先比较完整装载内容和调用计划，随后用去调试ELF作独立补充核对；三份原始ELF同时归档。\par


原ELF SHA256：\seqsplit{02159c7ef78cb31eb24cf706bf2974eadea313710563f235f2b11a7bb922de56}。\par
复制对照ELF SHA256：\seqsplit{cfad14278ee4a3590de9bcbe7589b808fd91959156cef93ba92a78a90db0d99b}。\par
SDK接入ELF SHA256：\seqsplit{f0ddf1167c961101e70f58cfd012ae788d2965324ad057aa4ade5af1de5bc42e}。\par

软件包接入按对象、链接和装载内容逐层核对。静态库只含完整gemv.o，保留16个函数、原节和重定位；链接映射与GNU依赖文件确认调用程序实际使用SDK头和库成员。这样可以直接定位最终函数来自哪里，以及重定位和符号解析如何完成。\par


最终比较包含原合格镜像、原源码复制对照和SDK接入版。三者BIN均为31440字节且逐字节一致，入口为0x80000000；PT\_LOAD、BSS、函数地址与机器码、48条调用表记录和12组输入数据也一致。装载内容相同，把接口改用SDK后的运行输入与原程序连接起来。\par


完整ELF的差异位于调试节。源文件路径、头文件位置和行号变化会改变这些信息，但装载段比较已经确认程序实际使用的字节相同。进一步用GNU objcopy去掉调试信息后，三份ELF也逐字节一致，形成与主检查不同路径的补充核对。\par

\subsection{接入检查器的正负控制}

7项负控全部被拒绝：实际ELF代码字节改坏、调用表改坏、启动入口改变、BSS大小改变、ISA属性改变、输入计划改变，以及真实GNU链接错误空库。另有只改调试字节的正控，确认不会将纯调试变化判为行为改变。这些控制样例用于验证主机端镜像检查器。\par


627项输入及封存运行材料绑定0f、EMU/NEMU与GNU工具。原33个矩阵进程和3个C进程的记录关系及输出树保持，接入过程另行保存GNU构建和最终镜像等价回执，形成软件包与既有功能结果之间的连接。\par


成功流水线：\seqsplit{2026-10-27T21:59:40.771873+08:00} 至 \seqsplit{2026-10-27T21:59:43.182911+08:00}，exit0。两次GNU构建，8份主流水线命令回执，另3份去调试GNU检查回执。流程首次通过，负控和原始文件全部保留。\par

\reportfigure{S3-08}{C8三方静态装载内容与调用输入对照}{17.9cm}{原合格夹具、同源重建与SDK实际链接版的BIN、PT\_LOAD和有序调用表、输入身份一致。PT\_LOAD按30024字节文件内容与1416字节零填充的真实比例绘制；零填充包含对齐尾部，不能全部称为BSS。调用及输入网格表示逻辑记录，不表示物理地址排列。调试节差异位于非装载区；627个受保护文件已核验。静态等价支持沿用既有夹具证据，SDK接入未新增EMU、NEMU或全核运行。}{精确SVG工程图；工程依据：原C8 relink RESULT及三方等价结构化记录；原图存档与所选资产身份见\seqsplit{figures/redesigned-20261105/selected-diagrams.json}。静态等价支持复用既有功能结果；不是新增全核运行}

\subsection{接入结果与功能记录复用}

C8包完成对象、可重定位链接和接回原C程序的完整镜像等价检查。原16函数、48例两轮96次正常调用的功能记录可据此复用，其中8×32占12例、24次调用。\par


该接入实验采用静态构建与镜像比较，EMU/NEMU运行次数为0。273 cycles与1.095238倍来自原第25条性能矩阵，描述原C8完整调用；这里的结果说明SDK接入保持了原测试程序的装载内容和调用方式。\par


软件包已完成从普通C入口到最终镜像的接入验证。面向新应用时，可固定新的输入布局、调用路径和最终镜像，使用同路径firstdot/RVV比较完整调用成本。现有其他尺寸继续采用原选择：4×8与16×64使用firstdot，4×13使用RVV。\par

\section{面积、时序与实现成本}

\subsection{评估条件与主要结果}

使用Nangate45 Typical 1.10V/25°C，Yosys0.9、ABC1.01及OpenSTA2.0.17。主档周期1ns、负载5fF，采用固定线负载模型，另有0.5ns、2ns及1fF、20fF敏感性档。\par


{\small\setlength{\tabcolsep}{3pt}\begin{longtable}{>{\raggedright\arraybackslash}p{0.27\linewidth}>{\raggedright\arraybackslash}p{0.3\linewidth}>{\raggedright\arraybackslash}p{0.35\linewidth}}
\caption{评估条件与主要结果}\label{tab:3-7}\\
\toprule
\textbf{指标} & \textbf{FU} & \textbf{完整VFEX4父级}\\ \midrule\endfirsthead
\caption[]{评估条件与主要结果（续）}\\
\toprule
\textbf{指标} & \textbf{FU} & \textbf{完整VFEX4父级}\\ \midrule\endhead
标准单元面积 μm² & 17,298.246 & 284,614.946\\
nominal r2r到达 ns & 1.163404 & 3.261218\\
主档i2r setup裕量 ns & −0.268520 & −2.393998\\
\bottomrule\end{longtable}}


父级面积包含FU及DIV。负裕量表明1ns参考目标仍有时序违例，整核Fmax、功耗和布线后时序为待测。H1相对O2独立FU面积降低2.33386\%，父级降低0.137349\%，同时保留部分路径退化。0f的取消关钟修复以正确性为依据，其周期不退化与成本变化分别留档。\par

\subsection{寄存状态与综合面积}

面积采用Nangate45 Typical库，在相同映射流程和约束下统计标准单元面积。VDot FU为17298.246平方微米，完整VFEX4父级为284614.946平方微米。父级包含FU、DIV和共享控制，因此父级数值描述集成成本，FU数值用于观察点积单元自身的结构变化。\par


0f相对修复前的父级面积减少264.936平方微米，约0.0930\%，FU面积不变。父级DFF数量增加一个，回溯到既有整数DIV的\seqsplit{Bypass\_final\_rem\_plus\_d\_reg}[2]在映射中的保留差异；两版对应DIV生成RTL相同。综合会跨逻辑进行优化和裁剪，因此面积变化由完整映射结果给出，并非简单按某段新增逻辑逐门相加。\par


H1的32位保存优化与0f的门控修复分别记录。前者直接利用固定结果布局减少状态，后者改变父级取消后的清理控制。评估它们时，既查看局部模块，也查看共享ready、门控及父级输出路径，能够观察局部优化在集成环境中的整体代价。\par

\subsection{时序约束与路径分类}

静态时序使用1ns周期、5fF负载作为名义档，并配置理想时钟、0.05ns时钟不确定度与transition、0.1ns输入输出delay、0.05ns输入transition和固定线负载模型。另设0.5ns、2ns以及1fF、20fF负载档，在同一映射网表上观察约束敏感性。\par


路径按输入、寄存器和输出分类，记为i2r、r2r、r2o、i2o，同时单列ready、门控使能及父级中的VDot子集。到达时间描述数据传播，要求时间描述路径预算，裕量是两者的差。比较时把路径类别、端点和约束放在一起，便于区分组合路径变化与外部预算变化。\par


名义档下，FU与父级所报i2r路径的setup裕量分别为−0.268520ns和−2.393998ns，当前映射结果尚未满足1ns目标。五档共70项类别与条件对照中，到达时间增大30项、减小15项、不变25项。条形图展示0f在1ns／5fF名义档下的14类时序裕量，负裕量位于零线左侧。五档完整70项数值、修复前后的到达时间差及实际起止点保留在图数据附件与完整对照表。\par

\reportfigure{S3-07}{Nangate45面积与时序裕量}{15.7cm}{Nangate45 Typical 1.10V、25°C下的布局前评估。条形图展示0f在1ns／5fF名义档的14类裕量，其中10类为负；五档70项完整数值保存在图数据附件。父级面积包含FU。}{\seqsplit{r3-div-gate/cost/results/summary.json}与timing.csv；源身份、SHA及完整数值见\seqsplit{figures/evidence-20261105/figure-data.json}。Nangate45 Typical 1.10V/25°C，布局前理想时钟/线负载；没有整核PPA、布线后时序、功耗或Fmax结果}

\subsection{父级关键路径与改进方向}

0f修复后的父级ready路径到达时间降低约47.603ps，而r2r、i2r及gate\_enable类别的最差到达分别增加约94.304ps、119.132ps和32.965ps。部分路径改善与部分路径退化同时存在，说明完整父级中的控制、算术和门控网络对修改有不同响应。\par


端点回溯显示，两版r2r最差路径均位于VFDiv，i2r最差起点由vsew输入转为src输入，gate\_enable最差起点则在不同VFDiv实例间切换。类别中的最差路径可能随版本改变，因此优化时应检查实际起止点、逻辑深度与扇出，而不是仅根据类别名称判断同一条路径缩短了多少。\par


这些布局前结果给出了下一步时序工作的具体入口。可以优先分析完整父级中限制名义档的路径，再以相同库和约束建立候选对照。布线寄生、多工艺角、整核Fmax和功耗属于后续物理评估内容；当前功能仿真、周期比较和静态时序分别描述已有工程的不同方面。\par

\subsection{取消关钟与成本结果}

DIV取消后的清理需要经历明确的状态转移。父级识别匹配取消，子单元获得完成清理所需的边沿，随后撤销在途状态并回到空闲关钟。局部测试在48个取消窗口及其后144个采样拍中检查多余钟脉冲，观察目标是取消处理之后的时钟活动是否按预期停止。\par


门控输入使能与沿后采样的子时钟掩码代表不同位置的信号。使能参与ClockGate控制，掩码描述采样时哪些子时钟处于活动状态。结合占用、取消和清理事件阅读这些记录，能够解释为什么取消当拍仍需要一个有效清理边沿，以及何时可以重新接收其他功能单元的请求。\par


脉冲观察回答控制行为问题，面积和STA回答实现代价问题。动态功耗还取决于电容、电压、频率及各节点开关活动，需要相应模型与活动输入。当前评估据已有记录给出关钟行为、周期保持和同约束映射结果，后续功耗实验可以在相同功能场景上继续采集活动，保持比较对象一致。\par

\subsection{门控修复的工程取舍}

0f解决了匹配DIV取消后的无效关钟保持。局部测试在48个取消窗口和其后144个采样拍中确认无多余钟脉冲，配套全核回归与周期比较进一步检查功能和执行代价。\par
修复的直接目标是恢复正确的取消与关钟行为。周期对照评估其执行代价，完整父级映射评估其面积与时序影响；动态功耗需要结合开关活动和功耗模型进一步测量。\par
最终工程组合采用0f修复：固定微基准周期保持一致，空闲关钟缺陷得到处理，父级面积小幅下降；名义档r2r、i2r及门控使能的到达时间分别增加约94ps、119ps和33ps。后续时序优化可从这些实际路径入手，在相同库与约束下建立新候选对照。生产硬件对应提交为\seqsplit{0f12df675378556670afd0dd2153f47fcb330610}。\par

\section{验证汇总与设计结论}

\subsection{功能、复现与性能检查汇总}

{\small\setlength{\tabcolsep}{3pt}\begin{longtable}{>{\raggedright\arraybackslash}p{0.27\linewidth}>{\raggedright\arraybackslash}p{0.3\linewidth}>{\raggedright\arraybackslash}p{0.35\linewidth}}
\caption{功能、复现与性能检查汇总}\label{tab:3-8}\\
\toprule
\textbf{项目} & \textbf{验证结果} & \textbf{方法与结论}\\\midrule\endfirsthead
\caption[]{功能、复现与性能检查汇总（续）}\\
\toprule
\textbf{项目} & \textbf{验证结果} & \textbf{方法与结论}\\\midrule\endhead
{源码与输入身份} & {通过} & {0f生产源码398项与clean副本及选定Git版本一致；完整工作树快照24157项覆盖源码和软件材料。}\\
{环境部署验证} & {通过，采用原记录} & {2026-10-09 hello/vadd日志重新核对；阶段一12文件SHA及220295598B波形记录完成定位。}\\
{指令规格} & {内部合同已实现} & {SEW8/VLEN128，8对signed int8点积，低32位结果、高96位清零；Q01–Q07保留为待官方澄清项。}\\
{参考模型} & {主机检查通过} & {重新检查独立模型及确定性数据，覆盖22组定向与10000组随机输入。}\\
{功能单元与父级} & {原记录复算通过} & {0f真实FU与VFEX4父级覆盖完整结果、反压、取消、身份、异常和严格DIV关钟；167份原载荷及4＋16负控记录纳入复核。}\\
{整核协同仿真} & {组合验收通过} & {13项原记录与1项新控制诊断，覆盖117个不同用例、124次执行、O2控制及负控恢复；宿主UTC问题按前述占用审计单列。}\\
{点积微基准} & {原周期记录复用} & {0f修复前后96对正式窗周期一致；独立batch128为221对678 cycles，加速比3.067873倍。}\\
{通用GEMV入口} & {输入复核与复现通过} & {四尺寸768窗口及96个普通C输出复现一致；33个矩阵进程与3个普通C进程构成原运行集合，4×13选择RVV。}\\
{8×32专用入口} & {最终镜像复核通过} & {627项输入材料与原普通C镜像绑定，BIN、PT\_LOAD、入口、48表项和12数据组一致；原性能矩阵为273对299 cycles。}\\
{独立位宽扩展} & {独立单元验证通过} & {s16x8/r64算术与包装扩展在原实验环境中复算；生产指令和NEMU采用INT8。}\\
{面积与时序} & {成本复核通过，时序待收敛} & {0f面积与五档STA的1008份文件回读；FU17298.246、父级284614.946平方微米，1ns名义档存在负裕量。}\\
{独立构建复现} & {同机复现通过} & {独立源码、依赖、EMU/NEMU、AM/klib和25个应用重建；复现使用同一主机，并复用工具及下载缓存。}\\
{entry-rollback} & {入口与回退明确} & {C8静态库入口、RVV回退、firstdot原入口及O1/B0历史对照均保留。}\\
\bottomrule\end{longtable}}

\reportfigure{S3-09}{稳定组合的验收范围与保留边界}{14.6cm}{实心圆表示有相应记录，包含未达标结果，不等于PASS；空心圆表示未测，横杠表示该软件证据行不评价物理实现。局部握手拍与整核、C8的DUT周期分别记录。全核组合采用13项既有记录与1项新诊断，C8接入为静态等价。旧UTC容量审计FAIL、墙钟稳定性未证实均保留；1ns名义档14类中10类负裕量，整核Fmax、布线后时序和功耗未测。}{精确SVG工程图；工程依据：R5验收矩阵、P5 RESULT和现有版本数据档案；原图存档与所选资产身份见\seqsplit{figures/redesigned-20261105/selected-diagrams.json}。证据复用与新执行分开；不暗示未采用候选进入稳定代码}

\Needspace{8\baselineskip}
\subsection{优化候选的实测取舍}

{\small\setlength{\tabcolsep}{3pt}\begin{longtable}{>{\raggedright\arraybackslash}p{0.27\linewidth}>{\raggedright\arraybackslash}p{0.3\linewidth}>{\raggedright\arraybackslash}p{0.35\linewidth}}
\caption{优化候选的实测取舍}\label{tab:3-9}\\
\toprule
\textbf{探索} & \textbf{结果} & \textbf{最终处理}\\ \midrule\endfirsthead
\caption[]{优化候选的实测取舍（续）}\\
\toprule
\textbf{探索} & \textbf{结果} & \textbf{最终处理}\\ \midrule\endhead
输入分离候选（P1/P1-R1） & 对照窗口偏差超过预设1\%限值 & 暂缓\\
单拍两槽候选（P4） & 47项成本检查中6项失败 & 未采用\\
单拍两槽修订（P4-R1） & 47项成本检查中10项失败 & 未采用\\
C8分段归约 & 实测273.75→273周期，一致性检查与3\%收益目标未通过 & 未采用\\
FP32到QKV探索（P3） & 未取得完整计时窗口，尚未确认完整调用能否在7200秒预算内完成 & 暂缓\\
\bottomrule\end{longtable}}


综合这些结果，工程选用0f＋firstdot V2＋原C8。应用评估包括点积微基准和固定GEMV调用；另有INT16独立单元扩展实验，生产指令保持INT8。\par

\subsection{阶段结论}

本项目完成自定义指令硬件接入、独立NEMU扩展、分层功能验证及等价RVV周期分析。O2提升独立吞吐，H1降低结果状态成本，0f处理DIV取消关钟，firstdot与C8改善固定软件调用。固定源、镜像和样本支持结果复算，1ns时序收敛及完整物理评估构成后续开发方向。\par


\href{https://github.com/cc77-builder/riscv/blob/eca41d7cd959f5f3c4f240e382919c492d1ccb96/team/patches/nemu/README.md}{NEMU补丁说明} · \href{https://github.com/cc77-builder/riscv/blob/eca41d7cd959f5f3c4f240e382919c492d1ccb96/team/spec/measurement.md}{测量合同} · \href{https://github.com/cc77-builder/riscv/blob/eca41d7cd959f5f3c4f240e382919c492d1ccb96/team/optimizations/o2/r3-div-gate/PERFORMANCE\_COST\_RESULT.md}{0f性能与成本} · \href{https://github.com/cc77-builder/riscv/blob/eca41d7cd959f5f3c4f240e382919c492d1ccb96/team/docs/version-comparison-r5-g5-20261027.md}{版本数据对照}\par

\section{工程复现资料}

\subsection{固定版本与主要运行产物}

测量软件以T13实现为基础，原文件及SHA-256保存在\seqsplit{team/references/zufan-t13/}，对应提交为\seqsplit{f090665130107282be0990a7664aebf26e491b87}。测量实验按自己的构建、镜像和运行记录建立样本集合。\par
该测量硬件源码提交为\seqsplit{4460300e9557ad20ac152fa67b1827c5d62802dd}，工作负载的构建提交另行记录。\par

NEMU以\seqsplit{a48a11f208ff81ad765591d83271d1b339fb3861}为基点，使用vdot补丁1.1.0及\seqsplit{xs-16ae9dd-vdot-runtime-tval-v1}平台配置。补丁绑定7个变化文件，SHA256为\seqsplit{6c165bf5f94ef0bfb4f894a5f06eb6d620d4a7e0e12f3b1400fa72811f5d24db}。\par

\begin{itemize}
\item 本章引用的历史资料发布提交为eca41d7cd959f5f3c4f240e382919c492d1ccb96。
\item 生产硬件提交为0f12df675378556670afd0dd2153f47fcb330610。
\item clean-02 EMU的SHA256为\seqsplit{acf672eb39f836da4000b253849e3214089e9c8e5198060a57d2cc6bbd393da9}。
\item reference的SHA256为\seqsplit{62b1707c52a4c323468481e0c1863b79f244d9a3fd54aeb7b7d2ecb5f01a35ca}。
\item C8最终BIN的SHA256为\seqsplit{56805506eb98b68269c73a3645d8f75dff73b697e30dec5f93ccc4716089ad34}。
\end{itemize}
以上标识对应本章采用的工程组合及主要运行产物，分项实验使用各自记录的构建与镜像身份。

\subsection{两槽结构对照的样本记录}

O1与O2使用同一组16份冻结镜像，共执行36项（32次首次镜像运行＋4次代表batch镜像重放），478个校验窗口，192条正式性能窗口。每场景每模式每方法为3数据集×4measured＝12条样本；first/warmup和重放完整保留，重放不增加独立样本数。\par


{\raggedright 周期降低按 $1 - \frac{\sum_i C_{\mathrm{O2,custom},i}}{\sum_i C_{\mathrm{O1,custom},i}}$，加速比按该模式本次实测的 $\frac{\sum_i C_{\mathrm{RVV},i}}{\sum_i C_{\mathrm{custom},i}}$。式中 $C$ 为窗口周期数，$i$ 为同一统计组内配对的正式窗口索引。每组逐\seqsplit{dataset/repeat/seq/}镜像配对，机器码、工作量及退休指令一致；未扣循环/屏障或空窗开销。所有均值由该批正式窗口计算。single对应完整调用，batch对应四路同输入、独立目的的批量点积。\par}


每个作业通过DiffTest、完整128位输出、非scratch寄存器、CSR、哨兵和计数窗口检查，并满足各自O1/O2协议及有限排空条件。manifest列出各内核使用的scratch集合，scratch\_control单独验证允许变化与受保护寄存器的划分。额外4次batch重放逐seq周期与退休数一致。30项主机测试通过；全部原始stdout/stderr、真实退出码、版本、输入和检查器摘要保留。\par

\subsection{硬件对照与工具身份}

{\small\setlength{\tabcolsep}{3pt}\begin{longtable}{>{\raggedright\arraybackslash}p{0.27\linewidth}>{\raggedright\arraybackslash}p{0.3\linewidth}>{\raggedright\arraybackslash}p{0.35\linewidth}}
\caption{硬件对照与工具身份}\label{tab:3-10}\\
\toprule
\textbf{项目} & \textbf{O1} & \textbf{O2}\\\midrule\endfirsthead
\caption[]{硬件对照与工具身份（续）}\\
\toprule
\textbf{项目} & \textbf{O1} & \textbf{O2}\\\midrule\endhead
{实测EMU SHA-256} & {\seqsplit{6699225b76ceae479060741026683bb03de069c1e10a4a94e463304ba1e322d0}} & {\seqsplit{ac0123280bf6c742d27a7650505bc1d91db350d1519806edfd0be9c7ea208a57}}\\
{构建源码提交} & {\seqsplit{a239dace720b32216e83fb5a514930d334c69a26}} & {\seqsplit{b218cbc6e6b135e4ff219a902a9284a59d3c82c9}}\\
{来源} & {从同一硬件输入重建} & {复用G3已验收固定EMU}\\
{\shortstack[l]{O1\_HANDOFF /\\O2\_ELASTIC}} & {1/0} & {0/1}\\
{\seqsplit{O1\_TRACE/O2\_TRACE}} & {1/0} & {0/1}\\
\bottomrule\end{longtable}}


429项硬件源码与21236项依赖逐字节相同，MinimalConfig/单核/VLEN128/单线程EMU；共同FUSED\_INPUT1、T12\_TRACE0、T13\_TRACE1。各自合法协议观察器只提供只读观察，共同T13提供全局周期/退休计数。\par


路径归一化后的实际Verilator/C++命令、七类工具版本日志一致；实际LLVM19.1.7编译器由构建过程/proc记录并核对。j2、12G JVM、O3/O0；最多两个单线程EMU。新O1采用外置产品目录，原G3的1687份非对象生成文件已恢复并验证原始SHA，旧O2 EMU不覆盖。构建记录分别列出硬件源码、实际工具和封存提交，便于对应实验输入与结果。\par

\subsection{离线复算与模型数据}

Windows原生Python 3.13.2全量离线复算通过，覆盖36项作业、478个窗口、192条正式性能窗口及4项重放。见原生复算回执。复算直接读取封存的作业和窗口记录，Linux构建来源与Windows复算回执相互对应。\par


交付时补齐了原冻结功能生成器依赖的两份模型语料，逐字节取自运行工具提交63b95a5并绑定Git blob及SHA-256；原镜像、测量日志和646份主封存payload均未改变。原30项主机测试及新增3项语料绑定检查共33项通过；首次封装复算失败和修复后成功记录均保留。Windows请用Python UTF-8模式执行：\par


\Needspace{3\baselineskip}
\begin{Verbatim}[breaklines,breakanywhere,breaksymbolleft={},breaksymbolright={},fontsize=\footnotesize,baselinestretch=1]
python -X utf8 -B performance/portable.py --archive evidence/20261019-performance-ab --campaign campaign-01 --output /absolute/new/portable-verification.json
\end{Verbatim}


上述相对路径适用于Windows交付目录；仓库内对应\seqsplit{team/optimizations/o2}下的路径。\par

\Needspace{8\baselineskip}
\subsection{取消关钟修复的对照输入}

{\small\setlength{\tabcolsep}{3pt}\begin{longtable}{>{\raggedright\arraybackslash}p{0.27\linewidth}>{\raggedright\arraybackslash}p{0.3\linewidth}>{\raggedright\arraybackslash}p{0.35\linewidth}}
\caption{取消关钟修复的对照输入}\label{tab:3-11}\\
\toprule
\textbf{对象} & \textbf{修复前} & \textbf{修复后}\\\midrule\endfirsthead
\caption[]{取消关钟修复的对照输入（续）}\\
\toprule
\textbf{对象} & \textbf{修复前} & \textbf{修复后}\\\midrule\endhead
{实现} & {O2两槽＋H1低32位保存} & {O2＋H1＋DIV取消关钟修复}\\
{硬件源码} & {\seqsplit{8b5024fcbd96986871cc40ce8c9fb71a67c495eb}} & {\seqsplit{0f12df675378556670afd0dd2153f47fcb330610}}\\
{EMU SHA256} & {\seqsplit{e37b3c267cae29036bda0b6952a17db2523ca0b3a573d45281341d4207df8747}} & {\seqsplit{2d1db77cc407486a2387a4a69da2498b68d44980239025eebcc02370c02826e8}}\\
\bottomrule\end{longtable}}


硬件429份Git源码、21236份依赖、每版1672份生成产物、实际编译命令和七类工具版本均经核验。生产源码仅ExeUnit.scala不同。\par
同一MinimalConfig/单核单线程/O2与T13观察配置、128MB、seed=1、正常NEMU和flash；208份冻结软件载荷不变。\par
正常NEMU SHA256：\seqsplit{efbed0774addc27eac38af1a19ab5d9c9dd5f179b5015f25ae7b80b966185bd4}。\par
生成物差异包括Git标识、只读诊断日志、顶层注释及标量FDiv公共子表达式内联；这些差异与生产源码变化一起纳入生成物审查。\par

\subsection{修复前后逐窗周期核对}

{\small\setlength{\tabcolsep}{3pt}\begin{longtable}{>{\raggedright\arraybackslash}p{0.15\linewidth}>{\raggedright\arraybackslash}p{0.18\linewidth}>{\raggedright\arraybackslash}p{0.18\linewidth}>{\raggedright\arraybackslash}p{0.18\linewidth}>{\raggedright\arraybackslash}p{0.18\linewidth}}
\caption{修复前后逐窗周期核对}\label{tab:3-12}\\
\toprule
\textbf{场景} & \textbf{修复前vdot均值} & \textbf{修复后vdot均值} & \textbf{两版RVV均值} & \textbf{两版RVV/vdot}\\\midrule\endfirsthead
\caption[]{修复前后逐窗周期核对（续）}\\
\toprule
\textbf{场景} & \textbf{修复前vdot均值} & \textbf{修复后vdot均值} & \textbf{两版RVV均值} & \textbf{两版RVV/vdot}\\\midrule\endhead
{单次点积1} & {71.000} & {71.000} & {85.250} & {1.200704×}\\
{依赖链256} & {1857.000} & {1857.000} & {5440.500} & {2.929725×}\\
{四路独立批量128} & {221.000} & {221.000} & {678.000} & {3.067873×}\\
{加载/计算/存储64} & {1127.500} & {1127.500} & {1133.000} & {1.004878×}\\
\bottomrule\end{longtable}}


单位均为DUT全局cycles。每场景三组数据×四轮，vdot和RVV分别配对；192正式窗口形成96对，每对差值均为0，退休指令与工作量相同。\par
共478个校验窗口，保留first/warmup；四次额外batch重放完全一致，不加入正式样本。36项实际EMU均通过完整128位结果、DiffTest、GOOD TRAP、窗口计数和O2生命周期/排空检查。\par
各场景加速比按 $\frac{\sum_i C_{\mathrm{RVV},i}}{\sum_i C_{\mathrm{vdot},i}}$ 计算，其中 $C$ 为窗口周期数，$i$ 为该场景中配对的正式窗口索引；保留共同计时开销。修复前后逐窗周期一致，表明0f在这组固定微基准中保持了原有性能。\par
详见逐窗差值、全部样本、复算回执。\par

\subsection{综合面积与完整接口}

修复前采用H1封存的205份生成、综合和STA文件，经逐字节核验后作为基线；修复后重新生成RTL并完成综合与STA，形成同条件对照。\par
两边使用相同完整生产接口、相同未修改O4脚本、\seqsplit{Yosys0.9/ABC1.01/OpenSTA2.0.17}及其动态库、Nangate45 typical 25°C/1.1V库。\par
ABC目标1000ps、BUF\_X1驱动、5fF映射负载。面积是Liberty单元面积之和，单位μm²，非芯片布局面积。\par


{\small\setlength{\tabcolsep}{3pt}\begin{longtable}{>{\raggedright\arraybackslash}p{0.12285714285714286\linewidth}>{\raggedright\arraybackslash}p{0.12285714285714286\linewidth}>{\raggedright\arraybackslash}p{0.12285714285714286\linewidth}>{\raggedright\arraybackslash}p{0.12285714285714286\linewidth}>{\raggedright\arraybackslash}p{0.12285714285714286\linewidth}>{\raggedright\arraybackslash}p{0.12285714285714286\linewidth}>{\raggedright\arraybackslash}p{0.12285714285714286\linewidth}}
\caption{综合面积与完整接口}\label{tab:3-13}\\
\toprule
\textbf{边界} & \textbf{修复前面积} & \textbf{修复后面积} & \textbf{增量} & \textbf{增量比例} & \textbf{DFF数量} & \textbf{ICG数量}\\\midrule\endfirsthead
\caption[]{综合面积与完整接口（续）}\\
\toprule
\textbf{边界} & \textbf{修复前面积} & \textbf{修复后面积} & \textbf{增量} & \textbf{增量比例} & \textbf{DFF数量} & \textbf{ICG数量}\\\midrule\endhead
{fu} & {17298.246} & {17298.246} & {+0.000} & {+0.000000\%} & {1742\allowbreak{}→\allowbreak{}1742} & {0\allowbreak{}→\allowbreak{}0}\\
{parent} & {284879.882} & {284614.946} & {-264.936} & {-0.092999\%} & {17171\allowbreak{}→\allowbreak{}17172} & {3\allowbreak{}→\allowbreak{}3}\\
\bottomrule\end{longtable}}


父级映射DFF增加1个，已回溯到既有整数DIV的\seqsplit{Bypass\_final\_rem\_plus\_d\_reg}[2]：旧网表该位无驱动和消费者，新网表保留DFF\_X1。两版对应DIV生成RTL相同，没有新增RTL寄存器声明或VDot状态位；该差异来自整体综合优化与裁剪，面积因此按完整映射结果统计。详情在归档dff-review.json。\par


两版VDot均保存32位结果，并在公开128位输出处补齐高96位零；算术与控制/perf寄存器分组保持一致。父级DIV使用自己的完整结果格式。\par
FU生成5份文件完全相同。父级58份生成文件仅ExeUnit.sv和\seqsplit{LogPerfEndpoint.sv}不同；后者只改变打印顺序，并被实际SYNTHESIS定义排除。\par

\subsection{五档时序详细结果}

{\raggedright 同一映射的五档条件为1ns/5fF、0.5ns/5fF、2ns/5fF、1ns/1fF、1ns/20fF。理想时钟、0.05ns uncertainty与transition、0.1ns输入/输出delay、0.05ns输入transition、5K\_hvratio\_1\_1线负载。parent testEnable=0，不额外排除\seqsplit{ready/reset/flush}路径。\par}
每版70条类别/档位路径对照，父级网表规范化经逐bit拓扑绑定；未知单元和动态未约束端点拒收。以下nominal到达时间包含相应外部输入延迟；表格按路径类别比较，各版实际最差端点另行列出。\par


{\small\setlength{\tabcolsep}{3pt}\begin{longtable}{>{\raggedright\arraybackslash}p{0.15\linewidth}>{\raggedright\arraybackslash}p{0.18\linewidth}>{\raggedright\arraybackslash}p{0.18\linewidth}>{\raggedright\arraybackslash}p{0.18\linewidth}>{\raggedright\arraybackslash}p{0.18\linewidth}}
\caption{五档时序详细结果}\label{tab:3-14}\\
\toprule
\textbf{边界/路径} & \textbf{修复前arrival(ns)} & \textbf{修复后arrival(ns)} & \textbf{差值(ps)} & \textbf{修复后slack(ns)}\\\midrule\endfirsthead
\caption[]{五档时序详细结果（续）}\\
\toprule
\textbf{边界/路径} & \textbf{修复前arrival(ns)} & \textbf{修复后arrival(ns)} & \textbf{差值(ps)} & \textbf{修复后slack(ns)}\\\midrule\endhead
{fu/i2o} & {0.636728} & {0.636728} & {+0.000} & {0.213272}\\
{fu/i2r} & {1.190323} & {1.190323} & {+0.000} & {-0.268520}\\
{fu/r2o} & {0.611317} & {0.611317} & {+0.000} & {0.238683}\\
{fu/r2r} & {1.163404} & {1.163404} & {+0.000} & {-0.240117}\\
{fu/ready} & {0.253743} & {0.253743} & {+0.000} & {0.596257}\\
{\seqsplit{parent/gate\_enable}} & {1.855397} & {1.888362} & {+32.965} & {-1.010883}\\
{parent/i2o} & {3.128214} & {3.164359} & {+36.145} & {-2.314359}\\
{parent/i2r} & {3.194971} & {3.314103} & {+119.132} & {-2.393998}\\
{parent/r2o} & {3.185769} & {3.187058} & {+1.289} & {-2.337058}\\
{parent/r2r} & {3.166914} & {3.261218} & {+94.304} & {-2.341112}\\
{parent/ready} & {0.295424} & {0.247821} & {-47.603} & {0.602179}\\
{parent/vdot\_i2r} & {1.511395} & {1.474949} & {-36.446} & {-0.549064}\\
{parent/\allowbreak{}vdot\_r2o} & {2.832115} & {2.776701} & {-55.414} & {-1.926701}\\
{parent/vdot\_r2r} & {1.345889} & {1.356106} & {+10.217} & {-0.433760}\\
\bottomrule\end{longtable}}


70条类别/档位比较中，到达时间变大30条、变小15条、相同25条。全部差值和端点见完整时序CSV。\par
nominal所报路径中，修复前有10类负slack，修复后有10类。两版在1ns目标下均未收敛；整核Fmax、布线后时序与功耗列为后续评估内容。\par


最差端点回溯已保存\seqsplit{path-endpoint-review.json}：前后r2r均是VFDiv的fsm\_q[0]到\seqsplit{quo\_m1\_iter\_q\_f64\_0}[1]，i2r最差起点从vsew输入变为src输入；gate\_enable最差起点从vfdivs\_1切到vfdivs\_0，终点仍是浮点DIV的ICG。这些端点属于完整父级的时序分析对象，端点变化与到达时间变化共同用于解释取舍。\par

\subsection{复算入口与实验归档}

主机检查共62项，包括周期配对30项、验收身份6项、成本审计20项和归档拒错6项，均通过。预检中曾发现验收JSON字段适配和生成范围遗漏日志模块两项问题，修订后完成复核；对应失败回执与成功结果一并保存。\par
原始周期根：\seqsplit{/home/riscv/riscv-lab/o2-r3-div-gate-performance-20261025}。\par
原始综合根：\seqsplit{/home/riscv/riscv-lab/o2-r3-div-gate-cost-20261025}。\par
原H1完整build：\seqsplit{/home/riscv/riscv-lab/o2-h1-build-before-r3-div-gate-20261024}。\par
原始stdout/stderr、真实退出码、命令、源/依赖/镜像/工具/库哈希及所有路径报告随归档保存；EMU/NEMU和EDA可执行文件保留在本地，用SHA绑定，不放进日常Git提交。\par
周期和成本复算入口分别为\seqsplit{performance/recheck.py}与cost/compare.py，使用独立输出路径保存新回执。\seqsplit{package\_evaluation.py} -{}-verify ARCHIVE -{}-receipt NEW\_JSON用于核对归档字节。完整归档保存源、工具、库、镜像和路径报告的对应关系，可据此复算统计或重新安排相同输入的实验。\par

\subsection{C接口与编译开关}

{\small\setlength{\tabcolsep}{3pt}\begin{longtable}{>{\raggedright\arraybackslash}p{0.29\linewidth}>{\raggedright\arraybackslash}p{0.65\linewidth}}
\caption{C接口与编译开关}\label{tab:3-15}\\
\toprule
\textbf{入口} & \textbf{目标}\\\midrule\endfirsthead
\caption[]{C接口与编译开关（续）}\\
\toprule
\textbf{入口} & \textbf{目标}\\\midrule\endhead
{xs\_c8\_gemv\_8x32，开关缺省或 0} & {\seqsplit{config\_reuse\_gemv\_8x32\_reuse\_vdot}}\\
{xs\_c8\_gemv\_8x32，\seqsplit{XS\_C8\_GEMV\_FORCE\_RVV=1}} & {\seqsplit{config\_reuse\_gemv\_8x32\_reuse\_rvv}}\\
{\seqsplit{xs\_c8\_gemv\_8x32\_rvv}} & {始终 \seqsplit{config\_reuse\_gemv\_8x32\_reuse\_rvv}}\\
\bottomrule\end{longtable}}


\seqsplit{XS\_C8\_GEMV\_FORCE\_RVV} 最终展开只能是字面量 0 或 1；空值、表达式和未知标识不受支持，编译时拒绝。独立 guard 为 \seqsplit{XS\_C8\_SELECTED\_GEMV\_H}，API 版本 \seqsplit{XS\_C8\_GEMV\_API\_VERSION=1}。\seqsplit{XS\_C8\_GEMV\_PROFILE}为静态配置标签，CPU能力由调用环境确认。修改宏或API版本后，相关翻译单元一同重编，以保持调用点与库选择一致。\par


C8使用独立的开关、guard和入口别名，与原\seqsplit{XS\_GEMV\_FORCE\_RVV}及xs\_gemv\_*接口并存。应用选择C8或firstdot时，分别使用对应的头文件与静态库。\par

\subsection{对象构建与接入命令}

独立 GNU 汇编参数固定为 -march=rv64gcv -mabi=lp64d -mno-relax -c。完整 gemv.o 为 27,160 字节，SHA256 为 \seqsplit{f8bc13daeead98a414f313600a5778c96ac95a66b5bcb5c714ea3b98e0d25ab3}。确定性 libxs\_gemv\_c8.a 只含一个 gemv.o，保留 16 个函数、原节和重定位；\seqsplit{config\_reuse\_gemv.h} 从合格 gemv.h 原字节复制，仅改产物文件名。\par


软件包绑定源基点 \seqsplit{c823a8786e052c6e26f6915563c42b6d730dd8e3}、选中 RTL \seqsplit{0f12df675378556670afd0dd2153f47fcb330610} 和 clean-02 配套身份。完整输入、工具身份及验收结果以 INPUTS.json、验证器输出和版本数据档案记录为准。验收同时检查对象内容、链接关系与运行输入身份。\par


包级静态检查覆盖完整对象、单库成员、缺省及0/1配置下两别名的GNU调用对象、可重定位链接等价、符号解析、错误配置拒绝和SDK搬移。example.c提供编译与链接示例。随后的小节给出接回原C程序的最终镜像等价结果，将包级检查与原矩阵、普通C调用的功能记录衔接起来。\par


接入新应用时，根据wrapper、分派、CSR保存、对象组合和最终布局的变化安排目标运行；性能比较同时计入调用者实际承担的工作。叶函数字节核对用于确认计算主体，完整镜像和调用计划则用于确认最终运行方式。\par

本包内设 \seqsplit{XS\_C8\_GEMV\_FORCE\_RVV=1} 或显式调用 \seqsplit{xs\_c8\_gemv\_8x32\_rvv} 后完整重编。撤销 C8 选用时，恢复原 firstdot 头和匹配库并完整重编。C8稳定别名对应8×32，其余尺寸沿用各自选定的firstdot或RVV入口。\par

\subsection{链接对象与检查入口}

\begin{itemize}
\item 原主程序只把底层头的包含改为 \seqsplit{selected\_gemv\_c8.h}。
\item 8×32 reuse vdot/RVV × 三数据集，共6个函数表项改用 xs\_c8\_gemv\_8x32 / \seqsplit{xs\_c8\_gemv\_8x32\_rvv}。其余42项、数据、期望值、顺序和校验不变，测试向量沿用原计划。
\item 真正从 libxs\_gemv\_c8.a(gemv.o) 链入完整合格对象；SDK程序源目录没有gemv.S，对象目录没有绕过库的gemv.o。GNU依赖文件和link map确认实际包含SDK头、实际抽取库成员。
\item 在全新输出中分别构建原源码复制对照和SDK接入程序，保留原AM/klib和链接次序；只有输出路径和记录链接来源的选项不同。
\end{itemize}


接入仅调整头文件与六个8×32函数表项，其余调用计划和数据保持一致。\par

\subsection{宿主运行记录与并发审计}

原14项全核记录的UTC容量检查出现80.119ms表观四路重叠，状态为\seqsplit{FAIL/WALL\_CLOCK\_LIMITED}。后续按既定合同重新执行\seqsplit{o2-control-diagnostic}，其余13项从原始日志重新解析，形成组合验收。\par


最终组合14项通过功能、NEMU差分、O2覆盖和负控后同镜像恢复。\textbf{组合采用13项复用记录与1项新执行记录}。选定记录的严格UTC并发峰值为3（上限3）；新单进程以同boot/PID的单调时钟包络验证，峰值1，包络1497.000秒。该耗时是宿主审计数据，不用于DUT加速比。\par


新控制诊断检查131072字节寄存器结果，退出0；独立离线审计再次通过。原13项的占用记录采用UTC，新控制诊断另有单调时钟包络。\par


本次UTC间隔与单调包络仍相差-62.888356秒（UTC更短）。因此墙钟稳定性仍未证实；新占用审计依赖同boot的单调包络，而非容差修正UTC。原UTC容量失败及两种时间基准的差异均保存在审计记录中。\par

\evidencefigure{SUP-S3-SUMMARY}{R5功能与成本验收摘要}{9cm}{功能条目记录局部与全核验证的采用情况；成本条目记录面积结果和1ns目标下的负裕量。}

\subsection{工程使用与回退资料}

当前1ns STA尚未收敛，功耗、整核Fmax和布线后时序待评估。INT16扩展已有独立单元结果，生产全核继续采用INT8。4×8和16×64的配置复用控制差异保留为后续分析方向，4×13采用RVV。内部规格Q01–Q07仍待官方澄清。\par


选定的0f＋firstdot V2＋原C8作为稳定工程组合，保留可使用入口与回退路径。后续候选围绕父级关键时序、软件调用配置及实际应用工作量开展比较，沿用已经明确的输入、镜像和样本组织。\par


冻结文件24157项，生产源码398项与0f及clean副本一致。快照在 \seqsplit{/home/riscv/riscv-lab/software-import-20261026/r5-g5-acceptance-01/frozen-workspace.tar.gz}，已逐项回读，含工作树文件但不含Git历史和外部依赖。使用/回退见README，逐项身份见数据记录，原始新增证据见manifest。失败尝试与原28条数据档案共同保存在工程记录中。\par
\subsection{技术资料索引}
{\small\noindent \href{https://github.com/cc77-builder/riscv/blob/347b1ae7eafb0f30c3b2cc1f8da8d9c8cdd3f98e/team/spec/measurement.md}{周期测量边界与统计合同}；\href{https://github.com/cc77-builder/riscv/blob/347b1ae7eafb0f30c3b2cc1f8da8d9c8cdd3f98e/team/optimizations/o2/performance/README.md}{O1与O2同镜像验证过程}；\href{https://github.com/cc77-builder/riscv/blob/347b1ae7eafb0f30c3b2cc1f8da8d9c8cdd3f98e/team/optimizations/o2/r3-div-gate/PERFORMANCE\_COST\_RESULT.md}{0f硬件性能与五档成本详细记录}；\href{https://github.com/cc77-builder/riscv/blob/347b1ae7eafb0f30c3b2cc1f8da8d9c8cdd3f98e/team/optimizations/o2/integration/firstdot-v2-0f/package/c8/CONTRACT.md}{C8调用ABI、对象及回退合同}；\href{https://github.com/cc77-builder/riscv/blob/347b1ae7eafb0f30c3b2cc1f8da8d9c8cdd3f98e/team/optimizations/o2/integration/firstdot-v2-0f/package/c8/relink/RESULT.md}{C8最终链接和镜像等价验证}；\href{https://github.com/cc77-builder/riscv/blob/347b1ae7eafb0f30c3b2cc1f8da8d9c8cdd3f98e/team/optimizations/o2/integration/firstdot-v2-0f/acceptance/RESULT.md}{R5组合验收和证据复用}。\par}
